\documentclass[a4paper,12pt]{article}
\usepackage[utf8]{inputenc}
\usepackage[T1]{fontenc}
\usepackage[polish]{babel}
\usepackage{graphicx}
\usepackage{amsmath}
\usepackage{amssymb}
\usepackage{mhchem}
\usepackage{hyperref}
\usepackage{fullpage}
\usepackage{footnote}

% Ustawienia marginesów - sensowne wartości
\setlength{\topmargin}{1.5cm}
\setlength{\oddsidemargin}{2cm}
\setlength{\evensidemargin}{2cm}
\setlength{\textwidth}{15cm}
\setlength{\textheight}{20cm}
\setlength{\headheight}{15pt}
\setlength{\headsep}{20pt}
\setlength{\footskip}{30pt}

\begin{document}

% Strona tytułowa
\begin{titlepage}
    \centering
    \vspace*{2cm}
    
    {\LARGE\bfseries Twój Tytuł Dokumentu\par}
    \vspace{1cm}
    
    {\large\bfseries\textit{Autor: Twoje Imię i Nazwisko\par}
    \textit{Przedmiot: Nazwa Przedmiotu\par}
    \textit{Data: 17 października 2026\par}}
    
    \vfill
    
    {\large\bfseries \href{https://przyklad.pl}{Przykładowa strona internetowa}\par}
    
    \vspace{2cm}
    
    {\small\bfseries \textbf{Przypis:} Ten dokument zawiera wszystkie wymagane elementy\par}
    
\end{titlepage}

\null
\thispagestyle{empty}

\section{Wprowadzenie}

Poniżej przedstawiono dokument zawierający wszystkie wymagane elementy.

\subsection{Ilustracja}

Używamy pakietu \texttt{graphicx} do wstawiania obrazów:

\begin{figure}[h!]
    \centering
    \includegraphics[width=0.8\textwidth]{przyklad.png}
    \caption{To jest przykładowy podpis pod obrazem}
    \label{fig:przyklad}
\end{figure}

Więcej informacji znajdziesz na stronie <\ref{fig:przyklad}>.

\subsection{Wzór matematyczny}

Poniższy wzór wykorzystuje polecenia \texttt{\label{}} i \texttt{\ref{}}:

\begin{equation}
    F = G \frac{m_1 m_2}{r^2}
\end{equation}

Zobacz wzór <\ref{eq:przyciazanie}>.

\subsection{Wzór chemiczny}

Używamy pakietu \texttt{mhchem} do zapisu wzorów chemicznych:

\[
    \ce{2H2 + O2 -> 2H2O}
\]

Inny przykład reakcji:

\[\ce{Zn + 2HCl -> ZnCl2 + H2}\]

\subsection{Odnośnik do strony internetowej}

Więcej informacji można znaleźć w dokumentacji LaTeX:

\begin{verbatim}
\url{https://www.latex-project.org}
\end{verbatim}

Albo bezpośrednio:

\url{https://www.overleaf.com}

\section{Dodatkowe informacje}

Poniżej jeszcze raz wszystkie elementy na liście:

\begin{itemize}
    \item \textbf{Strona tytułowa} - polecenia \texttt{\author{}}, \texttt{\title{}} i \texttt{\maketitle{}}
    \item \textbf{Ilustracja} - pakiet \texttt{graphicx} z odwołańami \texttt{\label{}} i \texttt{\ref{}}
    \item \textbf{Wzór matematyczny} - pakiet \texttt{amsmath}, \texttt{amssymb} z \texttt{\label{}} i \texttt{\ref{}}
    \item \textbf{Marginesy} - pakiet \texttt{fullpage} lub ręczne ustawienia
    \item \textbf{Odnośnik internetowy} - pakiet \texttt{hyperref}, polecenie \texttt{\url{}}
    \item \textbf{Przypis} - polecenie \texttt{\footnote{}}
    \item \textbf{Wzór chemiczny} - pakiet \texttt{mhchem}, polecenie \texttt{\ce{}}
\end{itemize}

\end{document}
